\section{Equivalence at one-loop order}
\label{sec:equiv}

As has been proven in Sec.~\ref{sec:ope}, the quasi-PDF and pseudo-PDF as well as their matching coefficients are related by a simple Fourier transform in Eq.~(\ref{eq:equiv}) and Eq.~(\ref{eq:ftc-C}). This relation is valid to all orders in perturbation theory. In this section we check the relations in Eq.~(\ref{eq:equiv}) and Eq.~(\ref{eq:ftc-C}) at one-loop order. We choose $\Gamma=\gamma^0$ for our main presentation, but also quote final results for the case $\Gamma=\gamma^z$.

In the Feynman gauge, we calculate the quark matrix elements of the unpolarized iso-vector quasi-PDF, pseudo-PDF, and light-cone PDF at one-loop order in dimensional
regularization with $d=4-2\epsilon$. The external quark state is chosen to be on-shell and massless, and we regularize the UV and collinear divergences by $1/\epsilon_{\mbox{\tiny UV}}\ (\epsilon_{\mbox{\tiny UV}}>0)$ and $1/\epsilon_{\mbox{\tiny IR}}\ (\epsilon_{\mbox{\tiny IR}}<0)$ respectively. The one-loop order Feynman diagrams are shown in Fig.~\ref{fig:diagram}.


\begin{widetext}

\begin{figure*}[t!]
	\centering
	\includegraphics[width=.9\textwidth]{images/quasi_diagram.pdf}
	\caption{One-loop Feynman diagrams for the quasi-PDF, \spcorr and pseudo-PDFs. The first one is named ``vertex", the second and third ones are named ``sail", and the last one ``tadpole". The standard quark self energy wavefunction is also included.}
	\label{fig:diagram}
\end{figure*}

In an on-shell quark state with momentum $p^\mu=(p^0=p^z,0,0,p^z)$, for $\Gamma=\gamma^0$, each diagram gives
\begin{align}  \label{eq:qtloopint1}
\tilde{Q}^{(1)}_\text{vertex}(\zeta,z^2,\epsilon)
&= {\mu^{2\epsilon} \iota^\epsilon \over 2p^0}\bar{u}(p)\int {d^d k\over (2\pi)^d}(-ig T^a\gamma^\mu) {i\over \slashed k} \gamma^0 {i\over\slashed k} (-i gT^a\gamma^\nu){ - i g_{\mu\nu} \over (p-k)^2}u(p)e^{-i k^z z}\ ,
  \\
\tilde{Q}^{(1)}_\text{sail}(\zeta,z^2,\epsilon)
&= {\mu^{2\epsilon} \iota^\epsilon\over 2p^0}\bar{u}(p)\int {d^d k\over (2\pi)^d} (ig T^a \gamma^0) {1\over i(p^z-k^z)} \left( e^{-ip^z z}- e^{-ik^z z}\right)  \delta^{\mu z} {i\over \slashed k} (-ig T^a \gamma^\nu) {-ig_{\mu\nu}\over (p-k)^2}u(p)
  \nonumber\\
&+ {\mu^{2\epsilon}\iota^\epsilon \over 2p^0}\bar{u}(p) \int {d^d k\over (2\pi)^d} (-ig T^a \gamma^\nu){i\over \slashed k} (ig T^a \gamma^0) {1\over i(p^z-k^z)} \left( e^{-ip^z z}- e^{-ik^z z}\right)  \delta^{\mu z}   {-ig_{\mu\nu}\over (p-k)^2}u(p)\ ,
  \nn \\
\tilde{Q}^{(1)}_\text{tadpole}(\zeta,z^2,\epsilon)
&= {\mu^{2\epsilon}\iota^\epsilon \over 2p^0}\bar{u}(p)\int {d^d k\over (2\pi)^d} (-g^2)C_F \gamma^0 \delta^{\mu z}\delta^{\nu z}\left({e^{-ip^z z}- e^{-ik^z z}\over (p^z-k^z)^2} - {ze^{-ip^z z}\over i(p^z-k^z)}\right){-ig_{\mu\nu}\over (p-k)^2}u(p)\ ,
  \nn
\end{align}
where $\iota = e^{\gamma_E}/(4\pi)$ is included to implement $\mu$ in the $\overline{\rm MS}$ scheme, $C_F=4/3$, and $T^a$ is the SU(3) color matrix in the fundamental representation. The second term in the brackets in the last line, which is proportional $z$, does not contribute to the loop integral as it is odd under the exchange of $p^z-k^z\to -(p^z-k^z)$. The quark self-energy correction is $\tilde{Q}_{\rm w.fn.}^{(1)}(\zeta,z^2,\epsilon) =\delta Z_\psi\: \tilde{Q}^{(0)}(\zeta,z^2)$ with the tree level matrix element $\tilde{Q}^{(0)}(\zeta,z^2)=e^{-i\zeta}$ and on-shell renormalization constant $\delta Z_\psi$,
\beq
\delta Z_\psi = {\alpha_sC_F\over 2\pi}\left(-{1\over2}\right)\left({1\over\epsilon_{\mbox{\tiny UV}}} - {1\over \epsilon_{\mbox{\tiny IR}}}\right)\ .
\eeq

After carrying out the loop integrals in Eq.~(\ref{eq:qtloopint1}) according to the method in Ref.~\cite{Ji:2017rah}, we obtain
\begin{align}\label{eq:1loopdiagram}
\tilde{Q}^{(1)}_{\rm vertex}(\zeta,z^2,\epsilon)
=\,& {\alpha_sC_F\over 2\pi}e^{\epsilon\gamma_E} \int_0^1 du\ (1-\epsilon)(1-u) e^{-iu \zeta}\Gamma(-\epsilon) 4^{-\epsilon} \big(\mu |z|\big)^{2\epsilon}
  \\
 =\,& {\alpha_sC_F\over 2\pi} e^{\epsilon\gamma_E}\biggl({\mu|z|\over2}\biggr)^{2\epsilon}
 \frac{(-1)\Gamma(2-\epsilon)}{\epsilon_{\mbox{\tiny IR}}}
 {1-i\zeta - e^{-i\zeta}\over \zeta^2}
  \,, \nonumber\\
\tilde{Q}^{(1)}_{\rm sail}(\zeta,z^2,\epsilon)
=\, & {\alpha_sC_F\over 2\pi}  e^{\epsilon\gamma_E}\left[(i\zeta) \int_0^1 du \int_0^1 dt (2-u) e^{-i(1-ut)\zeta} \Gamma(-\epsilon) 4^{-\epsilon}(t^2z^2\mu^2)^{\epsilon}\right.
  \nonumber\\
 &\left. -\int_0^1 du \ e^{-iu \zeta}\Gamma(-\epsilon) 4^{-\epsilon}(z^2\mu^2)^{\epsilon}  + \left({1\over\epsilon_{\mbox{\tiny UV}}} - {1\over \epsilon_{\mbox{\tiny IR}}}\right) e^{-i\zeta}\right]
  \nonumber\\
 =\,& {\alpha_sC_F\over 2\pi} e^{\epsilon\gamma_E} \left({\mu|z|\over2}\right)^{2\epsilon}
 \Bigg\{\frac{-\Gamma(1-\epsilon)}{\epsilon_{\mbox{\tiny IR}}}
 {2(1-\epsilon)\over 1-2\epsilon} \left[{1-e^{-i\zeta}\over -i\zeta} + e^{-i\zeta}(-i\zeta)^{-2\epsilon}\big(\Gamma(2\epsilon)-\Gamma(2\epsilon,-i\zeta)\big)
 \right]
  \nn\\
 & + {\Gamma(1-\epsilon)\over \epsilon_{\mbox{\tiny IR}} } \frac{e^{-i\zeta}}{\epsilon} + \left({1\over\epsilon_{\mbox{\tiny UV}}} - {1\over \epsilon_{\mbox{\tiny IR}}}\right) e^{-i\zeta} \Bigg\}
  \,, \nn\\
\tilde{Q}^{(1)}_{\rm tadpole}(\zeta,z^2,\epsilon)
=\, & {\alpha_sC_F\over 2\pi} e^{\epsilon\gamma_E}\left({\mu|z|\over2}\right)^{2\epsilon}  {\Gamma(1-\epsilon)\over \epsilon_{\mbox{\tiny UV}} (1-2\epsilon)}e^{-i\zeta}
 \ . \nn
\end{align}
For simplicity we have left out the tree level multiplicative spinor factor $\bar u\gamma^0 u$ when quoting one-loop results in \eq{1loopdiagram}, and will continue to do so for the \spcorr, quasi-PDF, and pseudo-PDF results quoted below.
Since $\alpha_s^{\rm bare}=\alpha_s(\mu) \mu^{2\epsilon} Z_g^2 = \alpha_s(\mu) \mu^{2\epsilon} + {\cal O}(\alpha_s^2)$ is $\mu$-independent we do not include $\mu$ as an argument in the bare functions.
In the final result for each term we have also specified whether $1/\epsilon$ factors (that remain after expanding about $\epsilon\to 0$) are IR or UV divergences. Combined with the wavefunction corection, the bare \spcorr $\tilde{Q}^{(1)}(\zeta,z^2,\epsilon)$ is
\begin{align} \label{eq:1loopioffe}
\tilde{Q}^{(1)}(\zeta,z^2,\epsilon)
&= {\alpha_sC_F\over 2\pi}e^{\epsilon\gamma_E} \left({\mu|z|\over2}\right)^{2\epsilon} \Bigg\{
 {3\over2}\left({1\over \epsilon_{\mbox{\tiny UV}}}-{1\over \epsilon_{\mbox{\tiny IR}}}\right)e^{-i\zeta}
 +  \frac{(-1)\Gamma(2-\epsilon)}{\epsilon_{\mbox{\tiny IR}}}
  \Biggl[ {(1-i\zeta - e^{-i\zeta})\over \zeta^2}
 \\
& \qquad\qquad\qquad\qquad\qquad\quad 
 +{2\over 1-2\epsilon}\left\{ {1-e^{-i\zeta}\over -i\zeta} + e^{-i\zeta}(-i\zeta)^{-2\epsilon}
 \big(\Gamma(2\epsilon)-\Gamma(2\epsilon,-i\zeta)\big) - \frac{e^{-i\zeta}}{2\epsilon} \right\} \Biggr]
 \Bigg\}\ . \nn
\end{align}
Note that as $\epsilon\to 0$ the terms in the innermost curly brackets have no $1/\epsilon$ term.  Also we can verify that in the local limit of the operator that the bare one-loop correction vanishes as expected by conservation of the vector current:
\begin{align} \label{eq:Qvcc}
\lim_{z\to 0} \tilde{Q}^{(1)}(z P^z,z^2,\epsilon)
  = \lim_{z\to 0}  {\alpha_sC_F\over 2\pi}e^{\epsilon\gamma_E} \left({\mu|z|\over2}\right)^{2\epsilon} \bigg\{
 {3\over2}\left({1\over \epsilon_{\mbox{\tiny UV}}}-{1\over \epsilon_{\mbox{\tiny IR}}}\right)
  + \frac{3}{2\epsilon_{\mbox{\tiny IR}}} \bigg\} = 0
 \,,
\end{align}
where we note that it is important that the $1/\epsilon_{\mbox{\tiny IR}}$ terms cancel since the assumption $\epsilon>0$ is only valid for the $1/\epsilon_{\mbox{\tiny UV}}$ term.
The corresponding bare pseudo-PDF is
\begin{align}\label{eq:1looppseudo}
{\cal P}^{(1)}(x,z^2,\epsilon)
 &= {\alpha_sC_F\over 2\pi} \left[-1+x + \frac{2}{(1-2\epsilon)} -
 \left( {2\over (1-2\epsilon)}{1\over (1-x)^{1-2\epsilon}}\right)^{[0,1]}_{+(1)}\right]e^{\epsilon\gamma_E}\left({\mu|z|\over2}\right)^{2\epsilon} \frac{\Gamma(2-\epsilon)}{\epsilon_{\mbox{\tiny IR}}} \theta(x)\theta(1-x)\\
 &\quad +{\alpha_sC_F\over 2\pi} \:
  e^{\epsilon\gamma_E}\left({\mu|z|\over2}\right)^{2\epsilon}{3\over2}\left({1\over \epsilon_{\mbox{\tiny UV}}}-{1\over \epsilon_{\mbox{\tiny IR}}}\right)\, \delta(1-x)
 \ .\nn
\end{align}
Since we will encounter plus functions over different domains below, we define a plus function at $x=x_0$ within a given domain $D$ so that
\beq
	\int_D dx\ \big[ g(x)\big]_{+(x_0)}^D\, h(x) = \int_D dx\ g(x) \left[ h(x) - h(x_0)\right]\,.
\eeq
(See \app{ep} for more details.)
It is straightforward to confirm that the bare pseudo-PDF satisfies the local vector current conservation, $\lim_{z\to 0} \int\! dx \, {\cal P}^{(1)}(x,z^2\mu^2,\epsilon) =0$, with the same cancellation as in \eq{Qvcc}.

Now, according to the relations between the quasi-PDF and the \spcorr or pseudo-PDFs in Eqs.~(\ref{eq:fac-tq-orig0},\ref{eq:equiv}), we can do a Fourier or double Fourier transform of the results in Eqs.~(\ref{eq:1loopioffe},\ref{eq:1looppseudo}) to get the quasi-PDF. Despite its straightforwardness, the Fourier transform is subtle and the details are provided in \app{ft}. Here we simply quote the result for the bare quasi-PDF,
\begin{align}  \label{eq:1loopquasi}
\tilde{q}^{(1)}(x,p^z,\epsilon)
	=& {\alpha_sC_F\over 2\pi}\left\{{3\over2}\left({1\over \epsilon_{\mbox{\tiny UV}}}-{1\over \epsilon_{\mbox{\tiny IR}}}\right)\delta(1-x) + {\Gamma(\epsilon+{1\over2})e^{\epsilon\gamma_E}\over \sqrt{\pi}}{\mu^{2\epsilon}\over p_z^{2\epsilon}}{1-\epsilon \over \epsilon_{\mbox{\tiny IR}} (1-2\epsilon)} \right.\\
	&\times\left.\left[|x|^{-1-2\epsilon}\left(1+x+{x\over2}(x-1+2\epsilon)\right)-  |1-x|^{-1-2\epsilon}\left(x+{1\over2}(1-x)^2\right) + I_3(x)\right] \right\}\,,\nn
\end{align}
where
\beq
	I_3(x) = \theta(x-1)\left(x^{-1-2\epsilon}\over x-1\right)^{[1,\infty]}_{+(1)} -\theta(x)\theta(1-x) \left(x^{-1-2\epsilon}\over 1-x\right)^{[0,1]}_{+(1)} - \delta(1-x) \pi\csc(2\pi\epsilon) + \theta(-x) {|x|^{-1-2\epsilon}\over x-1}\,.
\eeq
After some algebra one can confirm that the bare quasi-PDF satisfies local vector current conservation, with $\int\! dx\, \tilde{q}^{(1)}(x,p^z,\epsilon) = 0$. To verify this result one must carefully separate out $1/\epsilon_{\mbox{\tiny UV}}$ factors arising from requiring $\epsilon>0$ to obtain convergence at $x=\pm \infty$, and $1/\epsilon_{\mbox{\tiny IR}}$ factors that arise from requiring $\epsilon<0$ to obtain convergence at $x=1$.

An alternate method of obtaining the quasi-PDF is to directly calculate it from the Feynman diagrams by first Fourier transforming $z$ into $xp^z$. As a result, the factors $(e^{-ip^z z}- e^{-ik^z z})$ are transformed into $[\delta(p^z-xp^z)-\delta(k^z-xp^z)]$, and all the loop integrals reduce to $(d-1)$-dimensional ones. This is the procedure for the matching calculations of the quasi-PDF used in Refs.~\cite{Xiong:2013bka,Stewart:2017tvs,Wang:2017qyg}, and is distinct from doing the Fourier transformation after fully carrying out the integrals as in Eqs.~(\ref{eq:1loopdiagram}--\ref{eq:1loopquasi}). As a cross-check we have confirmed in \app{quasi} that we obtain the exact same bare quasi-PDF in \eq{1loopquasi} from both procedures.

Now we consider the $\epsilon$ expansion to obtain $\overline{\rm MS}$ renormalized results for the \spcorr, pseudo-PDF, and quasi-PDF. Expanding the \spcorr in $\epsilon$ we obtain
\begin{align}
 \tilde{Q}^{(1)}(\zeta,z^2,\epsilon)=\delta \tilde{Q}^{(1)}(\zeta,z^2,\mu,\epsilon_{\mbox{\tiny UV}})+\tilde{Q}^{(1)}(\zeta,z^2,\mu,\epsilon_{\mbox{\tiny IR}})+{\cal O}(\epsilon)
\end{align}
with the $\overline{\rm MS}$ counterterm and renormalized \spcorr given by
\begin{align} \label{eq:1loopiofferen}
\delta \tilde{Q}^{(1)}(\zeta,z^2,\mu,\epsilon_{\mbox{\tiny UV}})
&= {\alpha_sC_F\over 2\pi}
  e^{-i\zeta} {3\over2}  {1\over \epsilon_{\mbox{\tiny UV}}}
 \,, \\
\tilde{Q}^{(1)}(\zeta,z^2,\mu,\epsilon_{\mbox{\tiny IR}})
&= {\alpha_sC_F\over 2\pi} \Bigg\{
  \frac{3}{2} \Big(\ln{\mu^2 z^2 e^{2\gamma_E} \over 4} +1\Big) e^{-i\zeta}
  + \Bigl( -\frac{1}{ \epsilon_{\mbox{\tiny IR}} } - \ln\frac{z^2\mu^2e^{2\gamma_E}}{4} - 1 \Bigr) h(\zeta)
  + \frac{2(1 \!-\!i\zeta\!-\! e^{-i\zeta})}{\zeta^2}
  \nn\\
&\qquad\qquad\quad
  + 4 i \zeta e^{-i\zeta}\, {}_3F_3(1,1,1,2,2,2,i\zeta)
  \Bigg\}
  \,. \nn
\end{align}
Here ${}_3F_3$ is a hypergeometric function and the Fourier transform of $\big[(1+x^2)/(1-x)\big]_{+(1)}^{[0,1]}$ gives the function
\begin{align}
  h(\zeta) &=\frac{3}{2}\, e^{-i\zeta}
  + \frac{1+i\zeta-e^{-i\zeta}-2i\zeta e^{-i\zeta}}{\zeta^2}
  - 2 e^{-i\zeta} \big[ \Gamma(0,-i\zeta)\!+\!\gamma_E\!+\!\ln(-i\zeta) \big]
  \,.
\end{align}
For the position space PDF we have
\begin{align}
  Q^{(1)}(\zeta,\mu,\epsilon_{\mbox{\tiny IR}})
  &= -  {\alpha_sC_F\over 2\pi}  \frac{1}{  \epsilon_{\mbox{\tiny IR}} }  h(\zeta) \,.
\end{align}

Next we expand the bare pseudo-PDF from Eq.~(\ref{eq:1looppseudo}) in $\epsilon$ to obtain the $\overline{\rm MS}$ counterterm and renormalized pseudo-PDF as ${\cal P}^{(1)} (x,z^2,\epsilon)= \delta {\cal P}^{(1)} (x,z^2\mu^2,\epsilon_{\mbox{\tiny UV}})+{\cal P}^{(1)} (x,z^2\mu^2,\epsilon_{\mbox{\tiny IR}})+{\cal O}(\epsilon)$ with
\begin{align}\label{eq:1looppseudopdfb}
\delta {\cal P}^{(1)} &(x,z^2\mu^2,\epsilon_{\mbox{\tiny UV}})
	= {\alpha_sC_F\over 2\pi}\:
    {3\over 2\epsilon_{\mbox{\tiny UV}}} \: \delta(1-x)
  \,, \\
{\cal P}^{(1)} &(x,z^2\mu^2,\epsilon_{\mbox{\tiny IR}})
\nn\\
	=&{\alpha_sC_F\over 2\pi}\left\{\left({1+x^2\over 1-x}\right)^{[0,1]}_{+(1)}\left[- \left({1\over \epsilon_{\mbox{\tiny IR}}}+\ln{e^{2\gamma_E}\over 4}\right)-\ln(z^2 \mu^2)-1\right] - \left(4\ln(1-x)\over 1-x\right)^{[0,1]}_{+(1)}  +2(1-x)\right\}\theta(x)\theta(1-x)\nonumber\\
	&  +{\alpha_sC_F\over 2\pi}\left[{3\over2}\ln(z^2\mu^2) + {3\over2}\ln{e^{2\gamma_E}\over 4}+{3\over2}\right]\delta(1-x)
 \,.\nn
\end{align}
Note that the renormalized $\overline{\rm MS}$ pseudo-PDF depends explicitly on $\mu^2$, and satisfies the relation to the renormalized $\overline{\rm MS}$ \spcorr given in Eq.~(\ref{eq:pseudoRen}).
It is also interesting to note that having expanded in $\epsilon$, local vector current conservation is no longer satisfied by the limit of the renormalized $\overline{\rm MS}$ pseudo-PDF, since
\begin{align} \label{eq:Ppdfconsrv}
\lim_{z\to 0} \int\! dx \, {\cal P}^{(1)} (x,z^2\mu^2,\epsilon_{\mbox{\tiny IR}}) \simeq (3\alpha_s C_F/4\pi) \lim_{z\to 0} \ln(z^2\mu^2)
\end{align}
gives a divergent result. The same divergence is present in the one-loop $\overline{\rm MS}$ renormalized \spcorr in \eq{1loopiofferen}.  Although this is the case in $\overline{\rm MS}$, it does not need to be the case in other renormalization schemes.

For the quasi-PDF there are two methods that we can consider for the renormalized calculation, either expanding the bare result in \eq{1loopquasi} and renormalizing in $(x,p^z)$ space, or following our preferred definition in \eq{fac-tq-orig0} and Fourier transforming the renormalized \spcorr in \eq{1loopiofferen}.  Although these two approaches will lead to the same final result for $C$ for practical applications, there is a subtle difference that we will explain.

First consider the renormalization of the quasi-PDF done in $(x,p^z)$ space. Expanding Eq.~(\ref{eq:1loopquasi}) in $\epsilon$, and writing $\tilde{q}^{(1)}(x,p^z,\epsilon) =\delta \tilde{q}^{(1)}(x,\mu/|p^z|,\epsilon_{\mbox{\tiny UV}})+\tilde{q}^{(1)}(x,\mu/|p^z|,\epsilon_{\mbox{\tiny IR}}) +{\cal O}(\epsilon)$ allows us to identify the $\overline{\rm MS}$ counterterm and renormalized quasi-PDF as
\begin{align}\label{eq:qPDFren}
\delta \tilde{q}^{(1)}(x,\mu/|p^z|,\epsilon_{\mbox{\tiny UV}})
=&\,  {\alpha_sC_F\over 2\pi} {3\over 2\epsilon_{\mbox{\tiny UV}}}
 \left[\delta(1-x) - {1\over2}{1\over x^2}\delta^+\Big({1\over x}\Big) - {1\over2}{1\over (1-x)^2}\delta^+\Big({1\over 1-x}\Big) \right]
 \,, \\
\tilde{q}^{(1)}(x,\mu/|p^z|,\epsilon_{\mbox{\tiny IR}})
=&\, {\alpha_sC_F\over 2\pi}\left\{
\begin{array}{ll}
\displaystyle \left({1+x^2\over 1-x}\ln {x\over x-1} + 1 + {3\over 2x}\right)^{[1,\infty]}_{+(1)}- \left({3\over 2x}\right)^{[1,\infty]}_{+(\infty)} &\, x>1\\[10pt]
\displaystyle \left({1+x^2\over 1-x}\left[- {1\over \epsilon_{\mbox{\tiny IR}}} - \ln{\mu^2\over 4p_z^2} + \ln\big(x(1-x)\big)\right] - {x(1+x)\over 1-x}\right)^{[0,1]}_{+(1)}  &\, 0<x<1\\[10pt]
\displaystyle  \left(-{1+x^2\over 1-x}\ln {-x\over 1-x} - 1 + {3\over 2(1-x)}\right)^{[-\infty,0]}_{+(1)} - \left({3\over 2(1-x)}\right)^{[-\infty,0]}_{+(-\infty)} \quad &\, x<0
\end{array}\right.\nn\\[5pt]
& + {\alpha_sC_F\over 2\pi}\left[\delta(1-x) - {1\over2}{1\over x^2}\delta^+\Big({1\over x}\Big) - {1\over2}{1\over (1-x)^2}\delta^+\Big({1\over 1-x}\Big) \right] \left( {3\over2}\ln{\mu^2\over 4p_z^2} + {5\over2}\right)
 \,,\nn
\end{align}
The details of working out the $\epsilon$ expansion of \eq{1loopquasi} are provided in \app{ep}, including definitions of the plus functions and $\delta$-functions at $x_0=\pm \infty$ that appear in the result quoted here. The $\overline{\rm MS}$ quasi-PDF obtained in \eq{qPDFren} still satisfies vector current conservation
\begin{align} \label{eq:qpdfconsrv}
 \int\! dx\: \tilde{q}^{(1)}(x,\mu/|p^z|,\epsilon_{\mbox{\tiny IR}})  =0 \,.
\end{align}
This is obviously the case for the plus function terms which individually integrate to zero, and is also true for the combination of $\delta$-functions which appears in \eq{qPDFren}.

The renormalized $\overline{\rm MS}$ quasi-PDF in Eq.~(\ref{eq:qPDFren}) differs slightly from that obtained using our definition in Eq.~(\ref{eq:fac-tq-orig0}). Using Eq.~(\ref{eq:fac-tq-orig0}) and the renormalized \spcorr in \eq{1loopiofferen} we instead obtain
\begin{align}\label{eq:qPDFft}
\delta \tilde{q}'^{(1)}(x,\mu/|p^z|,\epsilon_{\mbox{\tiny UV}})
=&\,  {\alpha_sC_F\over 2\pi} {3\over 2\epsilon_{\mbox{\tiny UV}}}\delta(1-x)\,, \\
\tilde{q}'^{(1)}(x,\mu/|p^z|,\epsilon_{\mbox{\tiny IR}})
=&\, {\alpha_sC_F\over 2\pi}\left\{
\begin{array}{ll}
\displaystyle \left({1+x^2\over 1-x}\ln {x\over x-1} + 1 + {3\over 2x}\right)^{[1,\infty]}_{+(1)}- \left({3\over 2x}\right)^{[1,\infty]}_{+(\infty)} &\, x>1\\[10pt]
\displaystyle \left({1+x^2\over 1-x}\left[- {1\over \epsilon_{\mbox{\tiny IR}}} - \ln{\mu^2\over 4p_z^2} + \ln\big(x(1-x)\big)\right] - {x(1+x)\over 1-x}\right)^{[0,1]}_{+(1)}  &\, 0<x<1\\[10pt]
\displaystyle  \left(-{1+x^2\over 1-x}\ln {-x\over 1-x} - 1 + {3\over 2(1-x)}\right)^{[-\infty,0]}_{+(1)} - \left({3\over 2(1-x)}\right)^{[-\infty,0]}_{+(-\infty)} \quad &\, x<0
\end{array}\right.\nn\\[5pt]
& + {\alpha_sC_F\over 2\pi}\left[\delta(1-x)\left( {3\over2}\ln{\mu^2\over 4p_z^2} + {5\over2}\right) +{3\over2}\gamma_E\left({1\over(x-1)^2}\delta^+({1\over x-1}) + {1\over(1-x)^2}\delta^+({1\over 1-x})\right)\right]
 \nn .
\end{align}
To carry out this calculation we defined the Fourier transformation of the singular function $\ln(\zeta^2)$ as
\begin{align} \label{eq:subtlety}
\int {d\zeta\over 2\pi}\ e^{i x \zeta} \ln\zeta^2
= & \biggl[ \left.{d\over d\eta} \int {d\zeta\over 2\pi}\ e^{i x \zeta} (\zeta^2)^\eta \biggr] \right|_{\eta=0}
=\biggl[ \left.{d\over d\eta} {4^\eta\over \Gamma(-\eta)} {\Gamma(\eta+1/2)\over \sqrt{\pi}} {[\theta(x)+\theta(-x)]\over |x|^{1+2\eta}} \biggr]
\right|_{\eta=0}\\
=& \gamma_E \left[\left(-\delta(x) + {1\over x^2} \delta^+\Big({1\over x}\Big)\right) + \left(-\delta(x) + {1\over x^2} \delta^+\Big({1\over -x}\Big)\right)\right] \nn\\
& -\left[\left({1\over x}\right)_{+(0)}^{[0,1]} + \left({1\over x}\right)_{+(\infty)}^{[1,\infty]}\right]\theta(x) - \left[\left({1\over -x}\right)_{+(0)}^{[-1,0]} + \left({1\over -x}\right)_{+(\infty)}^{[-\infty,-1]}\right]\theta(-x)\,,\nn
\end{align}
where we have used the results in Eqs.~(\ref{eq:ft},\ref{eq:expansion}) to derive the second and last equalities, and took the limit $\eta\to 0^+$ or $\eta\to 0^-$ when needed.  This $\tilde{q}'^{(1)}(x,\mu/|p^z|,\epsilon_{\mbox{\tiny IR}})$ does not satisfy vector-current conservation, and is different from $\tilde{q}^{(1)}(x,\mu/|p^z|,\epsilon_{\mbox{\tiny IR}})$ in \eq{qPDFren} only by the $\delta$-functions at $x_0=\pm\infty$. Within the function domain $-\infty < x < \infty$, they are exactly the same.  We will see below that both \eq{qPDFren} and \eq{qPDFft} eventually lead to the same result for the one-loop matching coefficient.


The final ingredient we need for the matching calculations is the PDF, whose one-loop bare matrix element can be written as a sum of an $\overline{\rm MS}$ counterterm and renormalized matrix element, $q^{(1)}(x,\epsilon) = \delta q^{(1)}(x,\epsilon_{\mbox{\tiny UV}})+q^{(1)}(x,\epsilon_{\mbox{\tiny IR}})$, where
\begin{align}
 \delta q^{(1)}(x,\epsilon_{\mbox{\tiny UV}}) &= 	{\alpha_sC_F\over 2\pi} {1\over \epsilon_{\mbox{\tiny UV}}} \left({1+x^2\over 1-x}\right)^{[0,1]}_{+(1)}  \theta(x)\theta(1-x)
  \,, \\
 q^{(1)}(x,\epsilon_{\mbox{\tiny IR}}) &= 	{\alpha_sC_F\over 2\pi} {(-1)\over \epsilon_{\mbox{\tiny IR}}} \left({1+x^2\over 1-x}\right)^{[0,1]}_{+(1)} \theta(x)\theta(1-x)
  \,. \nn
\end{align}

With the above results in hand we can now determine the matching coefficients up to one-loop order.  Using \eq{ps-q-fact} we find
\beq \label{eq:pseudoC}
{\cal C}^{(1)}(\alpha,z^2\mu^2)= {\cal P}^{(1)}(\alpha,z^2\mu^2,\epsilon_{\mbox{\tiny IR}}) - q^{(1)}(\alpha,\epsilon_{\mbox{\tiny IR}}) \,.
\eeq
Therefore the matching coefficient relating the pseudo-PDF and PDF in the $\overline{\rm MS}$ scheme with $\Gamma=\gamma^0$ is\footnote{A one-loop analysis of the \spcorr in the coordinate space also recently appeared in Refs.~\cite{Radyushkin:2017lvu,Radyushkin:2018cvn}. Our factorization result for the \spcorr in \eq{io-Q-fact} has a similar form to the hard part of the reduced \spcorr found in Eq.(3.35) of Ref.~\cite{Radyushkin:2017lvu} and Eq.(17) of Ref.~\cite{Radyushkin:2018cvn}. It is therefore interesting to compare our ${\cal C}(\alpha,z^2\mu^2)/C_0(\mu^2z^2)$ and this hard part. Our $\overline{\rm MS}$ result  \eq{ps-c} differs from Refs.~\cite{Radyushkin:2017lvu} due to the presence of the $2(1-\alpha)$ term. The result in the final version of Ref.~\cite{Radyushkin:2018cvn} agrees with ours. \eq{ps-c} also agrees with the original result derived in Ref.~\cite{Ji:2017rah}, up to the addition of our $e^{2\gamma_E}$ terms. The result for ${\cal C}(\alpha,z^2\mu^2)$  in \eq{ps-c} should be used to extract an $\overline{\rm MS}$ PDF from an $\overline{\rm MS}$ result for the pseudo-PDF.}
\begin{align} \label{eq:ps-c}
{\cal C}(\alpha,z^2\mu^2)
	=& \left[ 1+ {\alpha_sC_F\over 2\pi}\left({3\over2}\ln(z^2\mu^2)+{3\over2}\ln{e^{2\gamma_E}\over 4}+{3\over2}\right)\right]\delta(1-\alpha)
 \\
 	&+ {\alpha_sC_F\over 2\pi}\left\{\left({1+\alpha^2\over 1-\alpha}\right)^{[0,1]}_{+(1)}\left[-\ln(z^2 \mu^2)-\ln{e^{2\gamma_E}\over 4}-1\right] - \left(4\ln(1-\alpha)\over 1-\alpha\right)^{[0,1]}_{+(1)}  +2(1-\alpha)\right\}\theta(\alpha)\theta(1-\alpha)
  \,. \nn
\end{align}
This result is independent of the infrared regulator as it must be.
We have also computed the matching coefficient for the $\Gamma=\gamma^z$ case, and it is ${\cal C}_{\gamma^z}(\alpha,z^2\mu^2) = {\cal C}(\alpha,z^2\mu^2) + \Delta {\cal C}_{\gamma^z}(\alpha,z^2\mu^2)$ with
\beq
\Delta {\cal C}_{\gamma^z}(\alpha,z^2\mu^2) = {\alpha_sC_F\over 2\pi} 2(1-\alpha)\theta(\alpha)\theta(1-\alpha)\,.
\eeq
Due to the $\ln(z^2\mu^2)\delta(1-\alpha)$ term in \eq{ps-c}, the matching coefficient for the $\overline{\rm MS}$ pseudo-PDF again displays the fact that there is not a smooth local limit as $z\to 0$.  It is possible to define a scheme other than $\overline{\rm MS}$ to ensure that this limit is smooth, reproducing a renormalization for $z\to 0$ that agrees with the fact that the local operator corresponds with a conserved current. One such scheme would be to simply multiply all $\overline{\rm MS}$ renormalization constants by $C_0(\mu^2 z^2)$, which would lead to a   \spcorr renormalized in a different scheme, and a corresponding different matching coefficient in \eq{ps-c} with a smooth $z\to 0$ limit. This is equivalent to studying the ratio of \eq{Qtratio} from the start as advocated in Ref.~\cite{Radyushkin:2017cyf,Radyushkin:2017lvu}. We will give explicit results for this scheme choice below. This modified scheme should not be confused with the strict definition of the $\overline{\rm MS}$ scheme.


From \eq{fac-tq} the corresponding relation for the matching coefficient for the quasi-PDF defined in Eq.~(\ref{eq:fac-tq-orig0}) is
\beq \label{eq:quasiC}
C^{(1)}(\xi,\mu/(|y|P^z))= \tilde q'^{(1)}(\xi,\mu/|y|P^z,\epsilon_{\mbox{\tiny IR}})-q^{(1)}(\xi,\epsilon_{\mbox{\tiny IR}})\,.
\eeq
Therefore using \eq{qPDFft} the matching coefficient relating the quasi-PDF and PDF is
\begin{align} \label{eq:quasi-c}
C\left(\xi, {\mu\over |y| P^z}\right)
 = &\, \delta\left(1-\xi\right)+{\alpha_sC_F\over 2\pi}\left\{
	\begin{array}{ll}
	\displaystyle \left({1+\xi^2\over 1-\xi}\ln {\xi\over \xi-1} + 1 + {3\over 2\xi}\right)^{[1,\infty]}_{+(1)}- \left({3\over 2\xi}\right)^{[1,\infty]}_{+(\infty)}
    &\, \xi>1
    \\[10pt]
	\displaystyle \left({1+\xi^2\over 1-\xi}\left[
     - \ln{\mu^2\over y^2P_z^2} + \ln\big(4\xi(1-\xi)\big)\right] - {\xi(1+\xi)\over 1-\xi}\right)^{[0,1]}_{+(1)}
    &\, 0<\xi<1
    \\[10pt]
	\displaystyle  \left(-{1+\xi^2\over 1-\xi}\ln {-\xi\over 1-\xi} - 1 + {3\over 2(1-\xi)}\right)^{[-\infty,0]}_{+(1)} - \left({3\over 2(1-\xi)}\right)^{[-\infty,0]}_{+(-\infty)} \quad
    &\, \xi<0
	\end{array}\right.\\[5pt]
	& + {\alpha_sC_F\over 2\pi}\left[\delta(1-\xi) \left( {3\over2}\ln{\mu^2\over 4y^2P_z^2} + {5\over2}\right) + {3\over2}\gamma_E\left({1\over(\xi-1)^2}\delta^+({1\over \xi-1}) + {1\over(1-\xi)^2}\delta^+({1\over 1-\xi})\right)\right]\,.\nn
\end{align}
Again this result is independent of the IR regulator as it must be. Here the plus function terms $\big[ g_1(\xi) \big]_{+(1)}^{[1,\infty]}$ and $\big[ g_2(\xi) \big]_{+(1)}^{[-\infty,0]}$ have integrands that converge for $\xi\to \pm \infty$, behaving as $g_i(\xi)\sim 1/\xi^2$.  Note that if we had instead used the renormalized $\overline{\rm MS}$ quasi-PDF calculated in Eq.~(\ref{eq:qPDFren}), we would obtain a different matching coefficient $C$ with different $\delta$-functions at $\xi=\pm\infty$. However, the $\delta$-functions do not contribute to the convolution integral in Eq.~(\ref{eq:fac-tq}) for any integrable PDFs. For example, to carry out the convolution with $1/\xi^2 \delta^+(1/\xi)$ we can use $\delta^+\big({1\over \xi}\big)  = \lim_{\beta\to0^+}  \delta\big({1\over \xi} - \beta\big)$, which
when plugged into the factorization formula gives
\begin{align}
&\lim_{\beta\to0^+}\int {dy\over |y|} {y^2\over x^2}
\, \delta\Big({y\over x} -\beta\Big) f_{u-d}(y)
=\lim_{\beta\to0^+}\beta f_{u-d}(\beta x)\,.
\end{align}
For the plus-function at $\infty$ using \eqs{plusinfinity}{plus1} we have
\begin{align}
\int_{-1}^{+1}\! {dy\over |y|} \left[{1\over (x/y)}\right]_{+(\infty)}^{[1,\infty]} f_{u-d}(y)
&= \lim_{\beta\to 0^+} \int_{-1}^{+1}\! {dy\over |y|}
\bigg[  {\theta(x/y-\beta) \over x/y}
+ {y^2\over x^2}\delta\Bigl({y\over x} -\beta\Bigr)\ln\beta  \bigg] f_{u-d}(y)
\\
&= \int_{-1}^{+1}\! {dy\over x} {y\over |y|} f_{u-d}(y)
+ \lim_{\beta\to 0^+}\beta f_{u-d}(\beta x)\ln\beta
\,. \nn
\end{align}
In the last line we dropped the $\theta(x/y-\beta)$ since at small $y$ our PDF behaves as $f_{u-d}(y)\sim y^{-1+a}$ with $0<a<1$. This also implies
\begin{align}
\lim_{\beta\to0}\beta f_{u-d}(\beta x) \propto\: & x^{-1+a}\lim_{\beta\to0} \beta^a  =0\,, \\
\lim_{\beta\to0}\beta f_{u-d}(\beta x)\ln\beta \propto\: & x^{-1+a}\lim_{\beta\to0} \beta^a \ln\beta =0\,,
\end{align}
which means that the distribution contributions evaluated at $\xi=\pm\infty$ in the matching coefficient $C$ give zero contribution.

Therefore, the matching coefficients calculated from the quasi-PDFs in Eq.~(\ref{eq:qPDFren}) and Eq.~(\ref{eq:qPDFft}) are the same in effect, and we can simply drop all the $\delta$-functions at $\xi=\pm\infty$ when plugging them into the factorization formula:
\begin{align} \label{eq:quasi-matching}
C^{\overline{\rm MS}}\left(\xi, {\mu\over |y| P^z}\right)
= &\, \delta\left(1-\xi\right)+{\alpha_sC_F\over 2\pi}\left\{
\begin{array}{ll}
\displaystyle \left({1+\xi^2\over 1-\xi}\ln {\xi\over \xi-1} + 1 + {3\over 2\xi}\right)^{[1,\infty]}_{+(1)}- {3\over 2\xi}
&\, \xi>1
\\[10pt]
\displaystyle \left({1+\xi^2\over 1-\xi}\left[
- \ln{\mu^2\over y^2P_z^2} + \ln\big(4\xi(1-\xi)\big)\right] - {\xi(1+\xi)\over 1-\xi}\right)^{[0,1]}_{+(1)}
&\, 0<\xi<1
\\[10pt]
\displaystyle  \left(-{1+\xi^2\over 1-\xi}\ln {-\xi\over 1-\xi} - 1 + {3\over 2(1-\xi)}\right)^{[-\infty,0]}_{+(1)} - {3\over 2(1-\xi)}\quad
&\, \xi<0
\end{array}\right.\\[5pt]
& + {\alpha_sC_F\over 2\pi}\delta(1-\xi) \left( {3\over2}\ln{\mu^2\over 4y^2P_z^2} + {5\over2}\right)\,.\nn
\end{align}
The use of \eq{quasi-matching} in the factorization formula is valid for any PDF that behaves as $\lim_{y\to 0} f(y,\mu) \sim y^{-1+a}$ with $a>0$.
We have also computed the matching coefficient for the $\Gamma=\gamma^z$ case, and it is given by $C_{\gamma^z}(\xi, \mu/(|y| P^z)) = C(\xi, \mu/(|y| P^z)) + \Delta C_{\gamma^z}(\xi, \mu/(|y| P^z))$ with
\begin{align}
\Delta C_{\gamma^z}(\xi, \mu/(|y| P^z))  = {\alpha_s C_F\over 2\pi} 2(1-\xi)
  \,\theta(\xi)\theta(1-\xi) \,.
\end{align}

Note that our result for the quark matching coefficient in $\overline{\rm MS}$ differs from that of Ref.~\cite{Wang:2017qyg} which is a pure plus function, but gives a convolution that does not converge, just as in the case of the quasi-PDF with a transverse momentum cutoff, see Ref.~\cite{Stewart:2017tvs}.

Since the renormalized pseudo-PDF and quasi-PDF satisfy the relation in Eq.~(\ref{eq:equiv}) by definition, $C\left(\xi, \mu/(|y| P^z)\right)$ and ${\cal C}(\alpha,z^2\mu^2)$ that are given by Eqs.~(\ref{eq:ps-c}, \ref{eq:quasi-c}) automatically satisfy the relation in Eq.~(\ref{eq:ftc-C}).

Besides, if one uses a scheme other than $\overline{\rm MS}$ for the quasi-PDF, such as the scheme obtained by absorbing $C_0$ into the $\overline{\rm MS}$ renormalization constant, then this will lead to a result for the matching coefficient that is a pure plus function and hence satisfies current conservation. Starting with \eq{quasi-c} and using \eq{C0} together with \eq{subtlety} we obtain
\begin{align} \label{eq:quasi-matching-alt}
C^{\rm ratio\!}\left(\xi, {\mu\over |y| P^z}\right)
= &\, \delta\left(1-\xi\right)+{\alpha_sC_F\over 2\pi}\left\{
\begin{array}{ll}
\displaystyle \left({1+\xi^2\over 1-\xi}\ln {\xi\over \xi-1} + 1 - {3\over 2(1-\xi)}\right)^{[1,\infty]}_{+(1)}
&\, \xi>1
\\[10pt]
\displaystyle \left({1+\xi^2\over 1-\xi}\left[
- \ln{\mu^2\over y^2P_z^2} + \ln\big(4\xi(1-\xi)\big)-1\right] +1+ {3\over 2(1-\xi)}\right)^{[0,1]}_{+(1)}
&\, 0<\xi<1
\\[10pt]
\displaystyle  \left(-{1+\xi^2\over 1-\xi}\ln {-\xi\over 1-\xi} - 1 + {3\over 2(1-\xi)}\right)^{[-\infty,0]}_{+(1)}\quad
&\, \xi<0
\end{array}\right.\,,
\end{align}
and for the $\Gamma=\gamma^z$ case,
\beq
\Delta C^{\rm ratio}_{\gamma^z}(\xi, \mu/(|y| P^z))  = {\alpha_s C_F\over 2\pi} \big[2(1-\xi)\big]^{[0,1]}_{+(1)} \,.
\eeq
While retaining current conservation in the renormalized quasi-PDF, \eq{quasi-matching-alt} can be used for example as input to the two-step matching procedure in the lattice calculation of PDF in Refs.~\cite{Alexandrou:2017huk}.
For the matching step, an equivalent procedure is to study the ratio given in \eq{Qtratio} in the $\overline{\rm MS}$ scheme from the start, as advocated in Ref.~\cite{Radyushkin:2017cyf,Radyushkin:2017lvu}, performing its matching onto the PDF, which will yield \eq{quasi-matching-alt}. This concludes our discussion of matching results and the equivalence between the quasi-PDF and pseudo-PDF at one-loop order.

\end{widetext}
